\documentclass{article}
\usepackage[T1]{fontenc}
\usepackage{lmodern}
\usepackage{graphicx}
\usepackage{amsmath,amssymb}
\usepackage[a4paper,margin=2cm]{geometry}
\usepackage[absolute,overlay]{textpos}
\setlength{\TPHorizModule}{1mm}
\setlength{\TPVertModule}{1mm}
\usepackage[indonesian]{babel}
\usepackage{hyperref}
\setlength{\emergencystretch}{2em}
% unit: Halaman 66

\begin{document}

% === Halaman 66 (python/native) ===

66

\% Nama file: LSB\_sisip.m

\% Penyisipan pesan di dalam citra dengan metode modifikasi LSB. Panjang

\% pesan (L) disimpan di dalam citra stegano. 

\% Asumsi: jumlah bit untuk merepresentasikan panjang pesan maksimal 30 bit.

\% Input: 1. Citra berwarna (RGB)

\%        2. Teks dari keyboard

\% Output: citra stego (harus disimpan dalam format BMP)

\% Diprogram oleh Rinaldi Munir, @Informatika STEI-ITB


clear all;

\% Baca citra cover dan tampilkan

namaFile = input('Ketikkan nama file citra cover: ', 's');

coverImage = imread(namaFile);

imshow(namaFile); title ('Citra cover ');


tinggi = size(coverImage, 1);

lebar = size(coverImage, 2);

total\_byte = tinggi * lebar * 3;

total\_bit\_embed = total\_byte;


\% Baca pesan

pesan = input('Ketikkan pesan: ', 's');


A. Program penyisipan pesan

\end{document}
